% !TEX root = ../main.tex
\section{Why Committed Trajectories Match Serial Execution}
\label{app:consistency}

This appendix spells out the argument behind the guarantee of \cref{sec:method:commit} and derives the savings stated in \cref{sec:method:overview}.

\paragraph{Conditions.}
The argument rests on four conditions.
First, decoding is deterministic: for a fixed input and configuration, the output that enters the history does not depend on which other requests are being decoded.
Second, the environment is deterministic: the result of a tool and the new environment depend only on the current environment and the call, and equal fingerprints mean equal environments.
Third, speculative executions touch only copies, and a copy becomes the environment only through adoption.
Fourth, nothing changes the committed history or environment between the final check and adoption.

\paragraph{Argument.}
We compare \method with serial execution window by window and assume that both start a window with the same history and environment.
The parent receives the same input in both executions and therefore produces the same output.
Its call is committed either by adopting the copy on which the dispatched call ran, which requires the two calls to be equal, or by executing the call directly.
In both cases the observation and the new environment are those of serial execution.
If the successor passes the check, it has the input and configuration of the serial turn $t{+}1$ and started from the same environment.
It therefore produces the same output, the same tool result, and the same new environment, and adoption makes exactly these the committed ones.
If the successor is discarded, it leaves nothing behind, and turn $t{+}1$ runs as in serial execution.
Either way, both executions end the window with the same history and environment, so the next window starts from the same point.

\paragraph{The parent's entry in the history.}
When a successor is committed, the history holds the parent's entry in the form the successor used, built from the dispatched call, instead of the parent's raw output.
The input check shows that both forms render to the same tokens.
Because the chat template renders each message from its own content and role, they also render identically in every later input.

\paragraph{If reasoning were kept in the history.}
The guarantee does not depend on the history policy of \cref{sec:motivation:input}.
If a harness kept reasoning in the history, the successor's input would never match and every successor would be discarded.
\method would remain lossless but would bring no gain, and it should be disabled.

\paragraph{Sampling.}
With stochastic decoding, the argument holds if the output of each turn is a deterministic function of its input, its configuration, and its random seed, and if the successor uses the seed of the turn it stands for.
A discarded successor must not be resampled or selected by its outcome, because this would change the distribution of committed outputs.
\method never does so.
All our experiments use greedy decoding.

\paragraph{Derivation of \cref{eq:saving}.}
Measure time from the start of the parent, and assume that the anchor is correct and ready before the parent ends, $\tau_t\le R_t$.
With speculative tool dispatch alone, the parent ends at $R_t$ and turn $t{+}1$ then takes $B_{t+1}$, for a total of $R_t+B_{t+1}$.
With \method, the successor starts at $\tau_t$ and ends at $\tau_t+B_{t+1}$, and the window ends when both turns are done, at $\max(R_t,\tau_t+B_{t+1})$.
The difference is $\min(B_{t+1},R_t-\tau_t)$.
A window cannot end before its longer turn, and without \method it takes at most twice as long as that turn, so the speedup of a window is at most $2\times$.
If the anchor is not ready before the parent ends, no successor is started and nothing is saved.

\paragraph{Contention.}
Now let each of two concurrent requests progress at rate $1/s$, and write $m=\min(B_{t+1},\sigma_t)$ with $\sigma_t$ and $B_{t+1}$ measured without contention.
If $B_{t+1}\ge\sigma_t$, the parent needs $s\sigma_t$ instead of $\sigma_t$ to finish after the anchor, and the successor completes work worth $\sigma_t$ in that time.
The window then ends $s\sigma_t+B_{t+1}-\sigma_t$ after the anchor instead of $\sigma_t+B_{t+1}$, which saves $(2-s)\,\sigma_t$.
If $B_{t+1}<\sigma_t$, the successor needs $sB_{t+1}$, the parent advances by $B_{t+1}$ in that time and then continues alone, and the window ends $\sigma_t+(s-1)B_{t+1}$ after the anchor, which saves $(2-s)\,B_{t+1}$.
In both cases a committed window saves $(2-s)\,m$.
A discarded successor contributes nothing and slows the parent for as long as both run.
If it would have run alone for $m'$ of the slack, the parent ends $(s-1)\,m'$ later.
With a commit rate $p$ and overlaps of similar length in both cases, a started window saves $p\,(2-s)\,m-(1-p)(s-1)\,m=(p-(s-1))\,m$ on average.
